\begin{abstract}
Modern GPU-centric systems run heterogeneous workloads, such as LLM inference, training, and vector search, whose performance and cost heavily depend on policies for memory placement, kernel scheduling, and observability. Today, these policies are hard-coded within framework-specific runtimes, proprietary vendor libraries, and monolithic kernel and device code, leaving operators unable to adapt them to changing workloads, SLAs, or hardware without invasive code changes and disruptions. Prior work demonstrated that externalizing CPU-side policies using eBPF significantly improves flexibility; however, directly applying this to GPUs is challenging. GPU driver mechanisms tightly couple resource management with low-level hardware interactions, and existing programmability approaches restricted to host-only or device-only execution fail to support coherent, end-to-end policies. We argue that the GPU driver/device layer is the only layer that simultaneously (i) has global, cross‑tenant visibility and (ii) controls all privileged mechanisms, and therefore it must become the narrow waist of the GPU policy plane.

We present \sys, a unified, safe, and efficient programmable policy runtime spanning both CPU hosts and GPU devices. \sys refactors the GPU driver to expose stable, verified eBPF hooks at key control points and offloads verified eBPF programs into GPU device via a lightweight, sandboxed PTX JIT runtime, coordinated through a shared cross-layer state abstraction. This approach enables workload-specific policies for memory offloading, kernel preemption, and fine-grained observability with safety guarantees and low overhead.

We implement \sys by extending the Linux driver and building a GPU-side runtime, and evaluate it across realistic workloads including LLM inference, GNN training, vector search, and multi-tenant scenarios. \sys-based policies improve throughput by up to $5\times$, reduce tail latency by up to $2\times$, and introduce only a few percent observability overhead. Our results demonstrate that unified CPU-GPU extensibility is a practical and essential foundation for future GPU-centric systems.
\end{abstract}

\section{Introduction}

GPUs now dominate the compute and capital costs of modern datacenters and edge servers, powering workloads that range from latency-sensitive LLM inference pipelines to large-scale training, graph analytics, and vector search. Across these workloads, end-to-end performance and efficiency are heavily determined not on raw FLOPs, but on \emph{policies}: how the system places and migrates data across device and host memory, how it admits and schedules kernels in multi-tenant environments, and how it collects observability signals for online adaptation. These policies must evolve continuously as models, SLAs, and hardware change. Yet today, they are largely fixed within proprietary vendor stacks and per-framework runtimes, making GPU infrastructure rigid and costly to operate.

In current GPU stacks, the host driver and device firmware implement fixed, opaque mechanisms for memory management and offloading, context scheduling, and performance event collection. Frameworks such as PyTorch or TensorFlow, as well as emerging inference runtimes like vLLM, build bespoke management logic atop these mechanisms, often defensively preallocating GPU memory or embedding scheduling heuristics tailored to a specific workload mix. Since underlying mechanisms expose no safe extension points, operators cannot implement workload-specific policies for memory placement or preemption precisely where decisions occur, such as during device memory migration or when low-priority tasks risk delaying latency-critical queries. Consequently, GPU deployments frequently suffer from unpredictable tail latency, underutilization, and operational disruptions whenever policies require updates or fine-tuning.

By contrast, modern CPU systems have steadily transitioned toward programmable policies. The extended Berkeley Packet Filter (eBPF) ecosystem supports dynamically loaded, verified programs for scheduling (\texttt{sched\_ext}), page-cache eviction (\texttt{cache\_ext}), TCP congestion control (\texttt{struct\_ops}), and a broad array of observability tools~\cite{gregg16,gregg2021computing,Bachl21,Zhong22,jia2023programmable}. These designs expose a small set of stable kernel mechanisms through structured attach points, enabling eBPF programs to implement policies dynamically, without requiring kernel reboots or code patches~\cite{schedext-docs,cacheext-docs}. Applying a similar concept to GPUs, treating the GPU driver as another kernel subsystem and attaching host-side eBPF programs to control its policies, appears promising. Conversely, GPU-only instrumentation frameworks such as CUPTI, NVBit, Neutrino, and eGPU~\cite{cupti,villa2019nvbit,neutrino,eGPU} provide rich telemetry via binary rewriting at the PTX or SASS level, yet primarily target profiling, lack unified safety guarantees, and do not integrate with host mechanisms.

Our key observation is that modern GPU appplications need an OS‑level narrow waist for policy, analogous to how the IP layer serves as the narrow waist of the network stack. The GPU driver and device execution environment together are the only layer with both global visibility across tenants and direct control over privileged mechanisms such as page migration, kernel admission, and preemption. We therefore treat GPU resource management as a first‑class kernel subsystem, and we make this layer programmable by exposing a small set of verified attach points and executing a unified eBPF policy IR across both the driver and GPU kernels over shared maps in unified memory.

We argue that the right abstraction for GPU-centric programmability is to treat the CPU host and GPU device as a single, coherent policy domain, supported by a unified, restricted eBPF intermediate representation, a single shared map abstraction in unified memory, and a common verifier and control plane. Policies implemented this way seamlessly span privileged driver mechanisms and in-kernel GPU execution, thus enabling true cross-layer programmability for GPU systems. We observe that eBPF's restricted instruction set, strong verifier guarantees, and rich ecosystem of existing policy tools make it a uniquely suited substrate for extending unified programmability from CPUs to GPUs.

However, extending host-only eBPF directly to GPU-centric systems fails for two fundamental reasons. First, critical decisions, such as thread-block scheduling, warp prioritization, or identifying pathological CUDA kernel memory behaviors, are made device-side within GPU kernels and on-chip firmware invisible to host-only instrumentation. Second, monolithic one-size-fits-all GPU drivers tightly integrate resource management with hardware-level interactions (e.g., MMU operations, interrupt handling, context management) and offer no stable, safe hooks; naively inserting probes risks GPU hangs, kernel panics, or subtle resource leaks.

Designing a general substrate for GPU policy programmability therefore introduces several challenges. \emph{C1: Safe and efficient driver extensibility.} We must retrofit mechanism-policy separation into large, performance-critical GPU drivers without destabilizing their low-level interactions with hardware. \emph{C2: Safe and efficient device-side execution.} Policy logic must execute directly within GPU kernels on critical memory access and synchronization paths, while ensuring bounded execution and memory safety in a massively parallel environment. \emph{C3: cross-layer state and abstractions.} Policies often require correlating host-side and device-side signals (e.g., page access patterns with warp-level stalls), necessitating a unified state abstraction such as time and virtual address, and efficient synchronization across the CPU-GPU boundary. \emph{C4: A single control plane.} The policy runtime must support unified verification, deployment, and debugging across both host and device execution contexts to avoid introducing additional fragmentation.

We present \sys, a unified, safe, and efficient policy runtime for GPU-centric systems designed explicitly to address these challenges. \sys exposes a small, stable set of GPU mechanisms for memory placement and migration, kernel admission and preemption, and kernel-level instrumentation as programmable attach points within an extended GPU driver. Policies are written as eBPF programs using a restricted instruction subset, statically verified for memory safety and bounded execution. The identical eBPF intermediate representation is then offloaded onto GPU devices: \sys translates verified bytecode into PTX, injects it via lightweight trampolines at kernel launch time, and executes it within a sandboxed GPU-side runtime. Both host-side and device-side policy programs share a common map abstraction based in unified memory, enabling low-latency, coherent state sharing without explicit data marshalling.

We implement \sys on Linux by extending \sysdriver using NVIDIA's open GPU kernel modules\cite{}, and constructing a GPU-side device code JIT and sandbox runtime. Driver modifications follow a \texttt{struct\_ops}-style approach, introducing carefully chosen attach points within memory management and scheduling paths without significant restructuring, delegating policy decisions to verified eBPF callbacks. On the GPU device, lightweight trampolines save minimal register state, invoke translated eBPF snippets, and restore state, introducing sub-microsecond overhead as measured in our microbenchmarks. Unlike prior GPU observability tools~\cite{cupti,villa2019nvbit,neutrino,eGPU} or userspace eBPF implementations, \sys provides a single unified policy abstraction across driver and device contexts, managed by one verifier and control plane.

We evaluate \sys across realistic GPU datacenter workloads including LLM inference pipelines, GNN training, vector search, and mixed-priority multi-tenant scenarios. Using \sys, we implement adaptive memory prefetching and eviction policies, fine-grained kernel preemption strategies, and bcc-style GPU kernel observability tools. \sys-based policies improve throughput by up to $5\times$ and reduce tail latency by up to $2\times$ compared to static heuristics, while instrumentation-only deployments incur only $2$–$3\%$ overhead. These results confirm that unified, eBPF-based policy programmability spanning CPU and GPU contexts represents a practical and powerful abstraction for managing GPU-centric datacenter systems.

\begin{itemize}[leftmargin=*,itemsep=0pt]
    \item We identify and characterize key programmability gaps in GPU stacks for memory management, scheduling, and observability, showing how current rigid implementations prevent adaptation to evolving workloads.
    \item We introduce \sys, the first unified eBPF-based policy runtime to span CPU and GPU contexts, that spans CPU and GPU execution contexts, combining verified driver hooks with safe, JIT-compiled execution of BPF programs inside GPU kernels over a shared state abstraction and single control plane.
    \item We demonstrate that \sys enables adaptive, workload-specific policies significantly improving performance, resource utilization, and tail latency on realistic LLM inference, training, and vector workloads, with low overhead and without requiring application or driver restarts.
\end{itemize}
